##############################################%#
               \section{Figurer}               %#
%##############################################%#

\subsection{Indsæt figur}

% Husk: filsti/filnavn.jpg
% Husk: Caption {Figurtext}
% Husk: Label - så du kan referere til billedet senere. Det anbefales at bruge {fig:label}, {tab:label}, {lst:label} - så er det nemmere at få fat i den rigtige. Kaldes med \ref{fig:meemoaldec}
% [here, top, bottom, page] - ! forsøger at tvinge latex til at placere her. Læs: https://www.overleaf.com/learn/latex/Inserting_Images#Positioning

\begin{figure}[htb!]                % Start figur
    \centering                      % Centrer
    \includegraphics[width=0.8\textwidth]{Figurer/Intro/insertimg.png} % Sæt størrelse, vælg fil
    \caption{Indsæt figur}          % Figurtekst
    \label{Intro/fig:insertfig}     % Label til referencer - brug sektion før label!
\end{figure}                        % Slut figur

\subsubsection{Indsæt to figurer side-by-side}

% To billeder ved siden af hinanden
\begin{figure}
    \centering
    \begin{minipage}{0.45\textwidth}
        \centering
        \includegraphics[width=0.9\textwidth]{example-image-a} % Venstre figur
        \caption{Venstre}
    \end{minipage}\hfill
    \begin{minipage}{0.45\textwidth}
        \centering
        \includegraphics[width=0.9\textwidth]{example-image-b} % Højre figur
        \caption{højre}
    \end{minipage}
\end{figure}

%##############################################%#
           \section{Code-snippets}             %#
%##############################################%#

% Indsæt kodesnippet fra fil: 
\lstinputlisting[language=c, firstline=50, lastline=74, firstnumber=50, caption=Server: sendFile(),label={lst:serversendFile}]{Filer/ex7/file_server.cpp}

% Indsæt kodesnippet direkte 
\begin{lstlisting}[caption={INDSÆT TEKST},captionpos=b]
--kode
\end{lstlisting}

%##############################################%#
               \section{Lister}                %#
%##############################################%#
% Lister
% Se https://www.overleaf.com/learn/latex/Lists
% indsæt \itemsep0em efter begin hvis du ikke ønsker stor linjeafstand mellem punkterne

% Liste i punktform 
\begin{itemize}
    \item Punkt
\end{itemize}

% Numereret liste
\begin{enumerate}
    \item Første ting her
\end{enumerate}

%##############################################%#
               \section{Math}                  %#
%##############################################%#

\begin{flalign}
f(u) & =\sum_{j=1}^{n} x_jf(u_j)&\\
     & =\sum_{j=1}^{n} x_j \sum_{i=1}^{m} a_{ij}v_i&\\
     & =\sum_{j=1}^{n} \sum_{i=1}^{m} a_{ij}x_jv_i
\end{flalign}

\begin{equation}
\frac{df}{dt}=\lim_{h \to 0}\frac{f(t+h)-f(t)}{h}
\label{math:test} 
\end{equation}

%##############################################%#
            \section{Overskrifter}             %#
%##############################################%#

% Overskrifter (der optræder i indholdsfortegnelsen):
\section{} % Nyt kapitel
\subsection{} % Underkapitel
\subsubsection{} % Under-underkapitel. 

%##############################################%#
              \section{Tabeller}               %#
%##############################################%#

% Tabeller - https://www.tablesgenerator.com/
% Ikke spild tid på at lave dem i hånden.

%##############################################%#
            \section{Kildehenvisning}          %#
%##############################################%#

% Kildehenvisning (biblatex)
% Se kildehenvisningsdokument for eksempler
\autocite{kildenavn}

% Print kildehenvisning
\renewcommand*{\mkbibnamefamily}[1]{\textsc{#1}}

\printbibliography[heading=bibintoc,title={Kildehenvisninger}] 

%Filtrer kildehenvisning
\printbibliography[heading=subbibintoc,type=article,title={Kun artikler}] %header=subbibintoc printer afsnittet i indholdsfortegnelsen

\printbibliography[type=book,title={Kun bøger}]

\printbibliography[keyword={datasheet},title={Datasheets}]
\printbibliography[keyword={latex},title={\LaTeX-relaterede ting}]

%##############################################%#
              \section{Infobox}                %#
%##############################################%#
% Infobox:
\begin{formal}
\texttt{\textbf{SOCK\_DGRAM}: Supports datagrams (connectionless, unreliable messages of a fixed maximum length).}
\end{formal}

